% Working paper draft. The venue style can be updated once a target is chosen.
\documentclass{article}
\usepackage[final]{iclr2025_conference}
\usepackage{times}

\usepackage{hyperref}
\usepackage{url}
\usepackage{booktabs}
\usepackage{graphicx}
\usepackage{amsmath,amssymb}
\usepackage{math_commands}
\usepackage{xcolor}
\usepackage{subcaption}
\usepackage{multirow}
\usepackage{xspace}

\newcommand{\method}{attention-local sinusoidal carrier\xspace}
\newcommand{\Method}{Attention-Local Sinusoidal Carrier\xspace}
\newcommand{\RMSNorm}{\operatorname{RMSNorm}}

\title{Attention-Local Sinusoidal Carriers Complement\\
Rotary Position Embeddings}

\author{%
Ethan Smith \\
\texttt{ethansmith2000@gmail.com}
}

\begin{document}

\maketitle

\begin{abstract}
Rotary position embeddings (RoPE) encode relative displacement by rotating
queries and keys, but their positional contribution remains multiplicatively
coupled to token content. We study a complementary additive path: at every
attention layer, a fixed sinusoidal carrier is added only to the inputs of the
query and key projections. The residual stream and values remain untouched,
the ordinary query and key matrices learn separate reads of the shared
carrier, and one learned scalar per layer controls its magnitude. The method
adds only one parameter per layer and preserves standard fused attention. In
a 153M-parameter causal language model trained for 3.28B tokens, it reduces
held-out negative log-likelihood by 0.0367 nats on average across three paired
training seeds. A rank-32 extension that learns a shared positional bottleneck
with separate query and key readouts improves by a further 0.0091 nats and
outperforms a parameter-matched feed-forward control. Factorial and location
ablations show that the carrier helps both with and without RoPE, but is most
effective when repeatedly exposed inside attention. Across seeds, the carrier
improves loss throughout the 1,024-token context and produces lower-entropy
attention with more mass on the first token and very recent keys. An exact
local logit decomposition shows substantial content--position cross terms,
rather than a purely position--position bias. These results support additive
absolute carriers and rotary relative geometry as complementary, though
partly overlapping, positional paths.
\end{abstract}

\section{Introduction}

Self-attention is permutation equivariant unless position enters its inputs or
logits. The original Transformer adds sinusoidal features to the residual
stream once at the model input \citep{vaswani2017attention}; many contemporary
language models instead rotate queries and keys using RoPE
\citep{su2024roformer}. RoPE is efficient and naturally expresses relative
offsets, but a position-only attention pattern must still be synthesized
through content-bearing query and key vectors. Analyses of trained RoPE models
suggest that different frequency bands acquire qualitatively different
positional and semantic roles \citep{barbero2025round}.

We ask whether RoPE benefits from a second, additive positional path that is
local to attention. Our construction repeatedly presents a fixed sinusoidal
carrier immediately before the query and key projections. It differs from
input positional embeddings because it never writes position into the
persistent residual stream, and it differs from a native head-space additive
carrier because the existing query and key matrices may independently mix its
frequency pairs. The construction is content independent, causal, compatible
with key--value caching, and requires no attention-logit tensor or custom
kernel.

The working contributions of this study are:
\begin{itemize}
    \item a minimal attention-local carrier requiring one learned scalar per
    layer and no changes to the attention kernel;
    \item paired evidence across seeds, model scale, and query-key
    normalization showing a persistent in-distribution language-modeling
    improvement over RoPE, plus a parameter-matched control for a stronger
    low-rank extension;
    \item controlled comparisons among residual-input, pre-projection, and
    head-space carrier locations, including the order of additive composition
    with RoPE; and
    \item three-seed mechanistic measurements of position-stratified loss,
    attention concentration, and an exact local decomposition of the
    content--position terms induced in attention logits.
\end{itemize}

\section{Background}

\subsection{Sinusoidal carriers}

For even model width $d$, define inverse frequencies
\begin{equation}
    \nu_i = \theta_s^{-2i/d}, \qquad i=0,\ldots,d/2-1,
\end{equation}
and the interleaved carrier
\begin{equation}
    \mathbf{s}(p)_{2i}=\cos(p\nu_i), \qquad
    \mathbf{s}(p)_{2i+1}=\sin(p\nu_i).
    \label{eq:carrier}
\end{equation}
We use the same base, $\theta_s=\theta_R=10{,}000$, as the RoPE backbone unless
stated otherwise. The two frequency banks are not identical: the carrier is sampled
across model width $d$, whereas RoPE is sampled across head width $d_h$. The
carrier is deterministic and depends only on the current position, so its use
cannot expose future token content.

\subsection{Rotary position embeddings}

For head width $d_h$, let $\omega_j=\theta_R^{-2j/d_h}$ and let $R(p)$ be the
block-diagonal rotation whose $j$th two-dimensional block has angle
$p\omega_j$. With the query-key normalization used in our primary backbone,
standard RoPE attention forms
\begin{equation}
    \mathbf{q}_p=R(p)\RMSNorm(W_q\mathbf{h}_p),\qquad
    \mathbf{k}_p=R(p)\RMSNorm(W_k\mathbf{h}_p),
    \label{eq:rope}
\end{equation}
so the query--key inner product depends on $R(p)^\top R(r)=R(r-p)$.
RoPE therefore exposes relative phase differences without materializing an
attention bias.

\section{Attention-Local Sinusoidal Carriers}

Let $\mathbf{h}^{(\ell)}_p$ denote the normalized residual input to attention
layer $\ell$. Our primary method uses
\begin{align}
    \mathbf{u}^{(\ell)}_p
        &= \mathbf{h}^{(\ell)}_p + \alpha_\ell\mathbf{s}(p), \\
    \mathbf{q}^{(\ell)}_p
        &= R(p)\RMSNorm\!\left(W_q^{(\ell)}
             \mathbf{u}^{(\ell)}_p\right), \\
    \mathbf{k}^{(\ell)}_p
        &= R(p)\RMSNorm\!\left(W_k^{(\ell)}
             \mathbf{u}^{(\ell)}_p\right), \\
    \mathbf{v}^{(\ell)}_p
        &= W_v^{(\ell)}\mathbf{h}^{(\ell)}_p.
    \label{eq:method}
\end{align}
Each $\alpha_\ell$ is a directly optimized, unconstrained fp32 scalar
initialized to one. It is shared between the query and key inputs, while
$W_q^{(\ell)}$ and $W_k^{(\ell)}$ provide separate learned reads. Thus an
$L$-layer model adds only $L$ parameters. Values and the persistent residual
state receive no positional addition.

Ignoring query-key normalization and writing $R_p=R(p)$, expanding one
attention logit yields four classes of terms:
\begin{align}
  &\left\langle R_pW_q(\mathbf{h}_p+\alpha\mathbf{s}_p),
          R_rW_k(\mathbf{h}_r+\alpha\mathbf{s}_r)\right\rangle \notag\\
  ={}&\langle R_pW_q\mathbf{h}_p,R_rW_k\mathbf{h}_r\rangle \\
   &+\alpha\langle R_pW_q\mathbf{h}_p,R_rW_k\mathbf{s}_r\rangle
    +\alpha\langle R_pW_q\mathbf{s}_p,R_rW_k\mathbf{h}_r\rangle \notag\\
   &+\alpha^2\langle R_pW_q\mathbf{s}_p,R_rW_k\mathbf{s}_r\rangle.
    \label{eq:fourterms}
\end{align}
The method can therefore express a position--position prior as well as two
content--position cross terms. Section~\ref{sec:mechanism} describes how we
measure their realized contributions with normalization restored.

\paragraph{Learned low-rank readout.}
Our enhanced variant retains the scalar path and adds a shared rank-$r$
positional bottleneck with separate all-head query and key readouts:
\begin{align}
    \widehat{\mathbf q}_p
      &= W_q(\mathbf h_p+\alpha\mathbf s_p)+U_qD\mathbf s_p, \\
    \widehat{\mathbf k}_p
      &= W_k(\mathbf h_p+\alpha\mathbf s_p)+U_kD\mathbf s_p, \\
    \mathbf q_p&=R(p)\RMSNorm(\widehat{\mathbf q}_p),\qquad
    \mathbf k_p=R(p)\RMSNorm(\widehat{\mathbf k}_p).
    \label{eq:rankreadout}
\end{align}
Here $D\in\mathbb R^{r\times d}$ is shared by the two branches and
$U_q,U_k\in\mathbb R^{d\times r}$ are separate within each layer. We use
$r=32$, initialize $U_q$ and $U_k$ to zero so the model begins exactly at the
scalar parent, and retain the same carrier and RoPE frequencies. This adds
589,824 parameters across eight layers (0.38\% of the model). The readout is
static and position-only; it is not a token-conditioned mapper.

\paragraph{A distinct head-space alternative.}
An additive carrier can instead be inserted after RoPE,
\begin{equation}
  \mathbf{q}_p=\RMSNorm\!\left(R(p)W_q\mathbf{h}_p+
  \mathbf{e}_q(p)\right),
  \label{eq:postrope}
\end{equation}
with an analogous key. This exposes a native head-space carrier but does not
pass it through $W_q$ or $W_k$. If the same indexed carrier is added before
RoPE, $R(p)\mathbf{e}(p)$ has doubled phase when its frequencies match RoPE;
the two orderings are consequently not interchangeable. This is not a pure
phase ablation when query-key RMS normalization has learned coordinatewise
gain $\Gamma$: although its scalar denominator commutes with rotation,
$R(p)\Gamma\neq\Gamma R(p)$ unless gains are equal within every rotary pair.

\section{Experimental Design}

The frozen protocol and run matrix are specified in
\texttt{paper/EXPERIMENT\_PLAN.md}. We summarize the main design here.

\subsection{Language-model architecture and optimization}

The controlled backbone is a decoder-only, pre-normalized Transformer with
multi-head causal attention, full-head RoPE, learned query-key RMS
normalization, GeGLU feed-forward blocks, and PyTorch fused scaled-dot-product
attention. The canonical model has width 768, eight layers and eight heads
(153.50M total parameters, of which 76.18M are outside the token embedding and
language-model head). A width-1024, 12-layer, eight-head model contains
305.63M parameters. Input and output token embeddings are untied.

We train on concatenated, 1024-token OpenWebText blocks using the GPT-2
tokenizer. The canonical evidence in Table~\ref{tab:mature} uses AdamW
learning rate $3\times10^{-4}$, betas $(0.9,0.98)$,
weight decay 0.01, 200 warmup steps, linear decay, global gradient norm 1.0,
and bf16 arithmetic. Our canonical cohort uses 32 sequences per update for
100k steps (3.277B nominal tokens), selected by a predeclared throughput and
memory benchmark. An earlier cohort uses eight sequences per update for 200k
steps (1.638B nominal tokens); we report it only as separate robustness
evidence because it differs in both token count and optimizer updates.

A subsequent bounded learning-rate audit kept every other setting fixed and
selected $1.2\times10^{-3}$ using RoPE development NLL rather than the carrier
contrast. We use that common recipe prospectively; we do not relabel the
completed $3\times10^{-4}$ cohort or claim an exhaustively tuned optimum.

The rank-32 readout uses a frozen $6.367487$ learning-rate multiplier on its
zero-initialized output factors, with inverse weight-decay compensation. This
value was chosen in a short, disjoint development screen to match the
function-space update scale of a wider rank-128 alternative. Rank 32 and rank
128 were subsequently tied at the mature endpoint, so the smaller variant was
frozen before the three-seed comparison.

The component ablation first completes a two-by-two factorial over carrier
presence and RoPE presence. Within the carrier-plus-RoPE method, it then tests
fixed unit magnitude versus learned magnitude, a globally shared versus
per-layer learned gate, and first-block-only versus repeated all-block
injection. Exploratory
learned phases, frequencies, and content-dependent carrier mappers are not
part of this primary matrix.

\subsection{Evaluation and uncertainty}

The primary endpoint is mean token negative log-likelihood on a frozen,
disjoint 1024-block holdout. Candidate and control use identical token order
and name-stable paired initialization. We report paired bootstrap intervals
over holdout blocks for within-run evaluation precision. Training seeds, not
individual blocks pooled across seeds, are the unit of replication. We also
report development curves, parameter count, peak memory, tokens per second,
and equal-wall-clock comparisons where throughput differs.

\section{Language-Model Results}

\begin{table}[t]
\centering
\caption{Canonical 153M-parameter results at 100k updates and batch 32. The
mean and sample standard deviation treat training seeds as replicates. Lower
NLL is better.}
\label{tab:mature}
\begin{tabular}{lrrrrr}
\toprule
Method & Seed 123 & Seed 456 & Seed 789 & Mean & SD \\
\midrule
RoPE & 3.20716 & 3.19940 & 3.21191 & 3.20615 & 0.00632 \\
Scalar carrier & 3.16955 & 3.16572 & 3.17323 & 3.16950 & 0.00375 \\
Matched FFN control & 3.16899 & 3.16606 & 3.16640 & 3.16715 & 0.00160 \\
Rank-32 carrier & \textbf{3.15941} & \textbf{3.15960} & \textbf{3.16218}
  & \textbf{3.16040} & 0.00155 \\
\bottomrule
\end{tabular}
\end{table}

The scalar carrier beats RoPE in all three seeds by $-0.03665$ nats on
average (descriptive seed-$t$ interval $[-0.04320,-0.03011]$). Rank 32 beats
the scalar carrier in all three seeds by $-0.00910$ on average
($[-0.01561,-0.00260]$), and beats the nearly parameter-matched FFN widening
in all three seeds by $-0.00675$ ($[-0.01344,-0.00007]$). Seed 123 selected
the low-rank candidate; the mean rank-32 advantage over scalar on only the two
subsequently registered seeds is $-0.00859$ nats.

The scalar method adds eight parameters and measured 213.7k target tokens/s
versus 216.3k for RoPE in the seed-123 run. Rank 32 adds 589,824 parameters
over scalar and measured 210.5k target tokens/s. Thus the minimal method has
negligible parameter cost and both variants retain the fused SDPA path; the
rank-32 extension has a modest parameter and throughput cost that should not
be obscured by the quality result.

\begin{table}[t]
\centering
\caption{Seed-123 learning-rate robustness at 100k updates. The final column
is a paired difference on the common selection window; all reported block-32
95\% intervals exclude zero. The $1.2\times10^{-3}$ recipe was selected using
RoPE development NLL.}
\label{tab:lrsensitivity}
\begin{tabular}{rrrr}
\toprule
Peak LR & RoPE NLL & Scalar NLL & Scalar $-$ RoPE \\
\midrule
$1.5\times10^{-4}$ & 3.26829 & 3.23912 & $-0.02916$ \\
$3.0\times10^{-4}$ & 3.19712 & 3.16076 & $-0.03636$ \\
$6.0\times10^{-4}$ & 3.16511 & 3.12033 & $-0.04478$ \\
$1.2\times10^{-3}$ & 3.15574 & 3.10127 & $-0.05447$ \\
\bottomrule
\end{tabular}
\end{table}

The carrier advantage therefore is not confined to its original learning
rate and grows across this bounded grid. Both methods improve through the
largest tested value, so the experiment establishes robustness and a
resource-bounded common recipe, not a bracketed optimizer optimum.

The earlier batch-8 cohort independently found a mean scalar improvement of
$-0.05533$ across three seeds. A one-seed 306M model retained a $-0.04058$
improvement, and a one-seed 153M ablation without query-key normalization
retained $-0.04952$. These are supportive transfer checks, not pooled
replicates of Table~\ref{tab:mature}.

\begin{table}[t]
\centering
\caption{Component and location ablations at the canonical 100k-step budget,
seed 123. The learned scalar carrier uses one gate per layer unless noted.}
\label{tab:components}
\begin{tabular}{llr}
\toprule
Carrier location & Positional backbone & Holdout NLL \\
\midrule
None & None & 3.26299 \\
None & RoPE & 3.18377 \\
Pre-Q/K scalar & None & 3.16979 \\
Pre-Q/K scalar & RoPE & \textbf{3.14676} \\
Pre-Q/K fixed unit gate & RoPE & 3.15322 \\
Pre-Q/K globally shared gate & RoPE & 3.14718 \\
Pre-Q/K first block only & RoPE & 3.14945 \\
Residual input, one shot & RoPE & 3.15339 \\
Fixed head-space, after projection & None & 3.16265 \\
Fixed head-space, after RoPE & RoPE & 3.15133 \\
\bottomrule
\end{tabular}
\end{table}

The two-by-two cells in Table~\ref{tab:components} show that the carrier helps
without RoPE and helps further when combined with RoPE. The interaction is
sub-additive: the joint NLL improvement is smaller than the sum of the two
individual improvements. Repeated pre-Q/K access is 0.00663 nats better than
writing the sinusoid once at the residual input. Learning the unit-initialized
gate helps by 0.00646, but using one global learned gate is within 0.00042 of
separate layerwise gates; the evidence therefore supports learned magnitude,
not the necessity of per-layer parameters. A fixed native head-space carrier
after RoPE is competitive but 0.00458 worse than pre-Q/K injection. This
comparison also changes carrier phase composition and the order of learned
query-key gains, as discussed around Equation~\ref{eq:postrope}.

\section{Mechanistic Analysis}
\label{sec:mechanism}

We evaluated all nine checkpoints from Table~\ref{tab:mature} on the same
1,024-block frozen holdout. Attention statistics use 64 evenly spaced blocks
chosen before inspection. The training seed is the replication unit; blocks,
layers, and heads are not treated as independent runs.

\begin{table}[t]
\centering
\caption{Seed-mean NLL differences by next-token target position. Negative
values favor the first method named.}
\label{tab:positionloss}
\begin{tabular}{lrrr}
\toprule
Target position & Scalar $-$ RoPE & Rank 32 $-$ RoPE & Rank 32 $-$ scalar \\
\midrule
1--15    & $-0.02066$ & $-0.02923$ & $-0.00857$ \\
16--31   & $-0.03792$ & $-0.04181$ & $-0.00388$ \\
32--63   & $-0.03234$ & $-0.04144$ & $-0.00910$ \\
64--127  & $-0.03489$ & $-0.04672$ & $-0.01183$ \\
128--255 & $-0.03582$ & $-0.04406$ & $-0.00824$ \\
256--511 & $-0.03625$ & $-0.04604$ & $-0.00979$ \\
512--1023& $-0.03795$ & $-0.04679$ & $-0.00883$ \\
\bottomrule
\end{tabular}
\end{table}

The scalar improvement is present in every position bin and each training
seed; it is smaller for the first 15 targets, where little context is
available. The rank-32 extension improves the seed mean in every bin. It is
therefore not obtaining its endpoint gain only from late-context positions.

\begin{table}[t]
\centering
\caption{Attention geometry averaged first over sampled blocks, layers, and
heads within a seed, then over three seeds.}
\label{tab:attentiongeometry}
\begin{tabular}{lrrr}
\toprule
Metric & RoPE & Scalar & Rank 32 \\
\midrule
Normalized entropy & 0.61563 & 0.57592 & 0.57657 \\
First-token mass & 0.01517 & 0.02731 & 0.02979 \\
Same-token mass & 0.06483 & 0.07083 & 0.07059 \\
Distance 1--3 mass & 0.19829 & 0.22227 & 0.22254 \\
Distance 64--255 mass & 0.24340 & 0.22151 & 0.22153 \\
Distance 256+ mass & 0.11370 & 0.10877 & 0.10903 \\
\bottomrule
\end{tabular}
\end{table}

Relative to RoPE, both carrier models have lower normalized attention entropy
in every seed, more first-token mass in every seed, and consistently move mass
from distances 64--255 toward distances 1--3. We do not summarize this as a
universal shortening of attention: expected attended distance increases for
one seed and decreases for two. Rank 32 changes these coarse attention
statistics little relative to scalar despite its additional NLL gain.

To restore query-key normalization in Equation~\ref{eq:fourterms}, we split
each raw projected vector into content and positional components, divide both
by the common RMS denominator computed from their trained sum, apply the
learned RMS gain and the same RoPE rotation to each, and then form
$L_{cc},L_{cp},L_{pc},L_{pp}$. The sum is algebraically exact and matches
logits recomputed directly from full Q/K to maximum absolute error
$4.01\times10^{-5}$. Table~\ref{tab:logitdecomposition} reports RMS after centering
each query row over its visible keys, which removes score shifts that cannot
affect softmax.

\begin{table}[t]
\centering
\caption{Exact local logit decomposition, averaged over three seeds. ``Pos.
total'' is $L_{cp}+L_{pc}+L_{pp}$; KL compares full local attention with
$\operatorname{softmax}(L_{cc})$.}
\label{tab:logitdecomposition}
\begin{tabular}{lrr}
\toprule
Metric & Scalar & Rank 32 \\
\midrule
$L_{cc}$ centered RMS & 1.3441 & 0.9131 \\
$L_{cp}$ centered RMS & 0.3444 & 0.5500 \\
$L_{pc}$ centered RMS & 0.4512 & 1.0961 \\
$L_{pp}$ centered RMS & 0.3187 & 1.1476 \\
Pos. total centered RMS & 0.8366 & 1.6449 \\
Cosine(pos. total, full logits) & 0.8234 & 0.8638 \\
KL(full $\Vert$ content-only) & 0.7059 & 1.9133 \\
\bottomrule
\end{tabular}
\end{table}

Both methods make substantial use of the two content--position cross terms;
the carrier is not well described as only adding a content-independent
relative bias. Rank 32 moves considerably more local score variation into the
explicit positional and cross terms, while reducing the content--content RMS.
This is a descriptive decomposition conditioned on each trained hidden state,
not an independently runnable counterfactual model.

We therefore pair it with checkpoint interventions. Across three seeds,
replacing the rank-32 readout by its positionwise mean costs only 0.01109 NLL,
whereas removing that mean costs 3.47553 and zeroing the complete direct path
costs 3.51250. Zeroing the trained scalar endpoint changes NLL by only
0.00011. Yet a separately trained rank-32 model without the scalar anchor is
0.00830 worse than the full method at seed 123 and trails throughout training.
Together these results suggest that the scalar path is useful optimization
scaffolding, while the learned readout is dominated by a pre-RoPE constant
component whose rotation still creates a relative-position kernel. Ordinary
learned Q/K biases do not substitute for it: adding them to the scalar parent
changes NLL by only $+0.00010$ in the matched seed-123 control.

Finally, we tested whether the retained carrier solutions depend narrowly on
the exact sinusoidal origin. Replacing $\mathbf s(p)$ by
$\mathbf s(p+c)$ costs at most 0.00044 mean NLL for $c\in\{1,4\}$ across both
methods, and both retain most of their RoPE advantage at $c=64$. Much larger
out-of-distribution shifts are damaging: the scalar model crosses above RoPE
around $c=1024$, while rank 32 crosses around $c=256$. This is checkpoint
sensitivity rather than randomized-origin training evidence; it rules out
fragility to a neighboring phase origin, not dependence on the broad absolute
phase region observed during training.

\section{Related Work}

Absolute sinusoidal embeddings originate with the Transformer
\citep{vaswani2017attention}, while RoPE encodes relative offsets through
query-key rotations \citep{su2024roformer}. Recent analyses study how RoPE
frequency bands are actually used and motivate partial rotation
\citep{barbero2025round}. Other work learns context-dependent positions
\citep{golovneva2024cope}, mixes RoPE and position-free attention across
layers \citep{yang2025ropenope}, or adapts RoPE to two-dimensional vision
tokens \citep{heo2024ropevit}. Our focus is complementary but functionally
overlapping: a fixed absolute carrier is repeatedly made available inside
each attention block while the standard RoPE rotation remains intact. The
camera-ready literature review will additionally cover additive attention
decompositions, learned Fourier features, relative attention biases, learned
RoPE spectra, and SDPA-compatible Fourier priors.

\section{Limitations}

The replicated evidence uses one corpus, one fixed training context, one
153M-parameter architecture, and autoregressive language modeling. Scale and
query-key-normalization transfer each use one seed under the earlier batch-8
recipe. Component, location, and optimization-path ablations also use one
seed. The selected rank-32 method has only two fresh confirmatory seeds because
seed 123 participated in design selection. Evaluation-block bootstrap
intervals do not measure training-seed variation, and three-seed $t$ intervals
remain low-powered. The models are controlled research Transformers rather
than replicas of a production LLM. We therefore make no claim of universal
positional-encoding superiority, length extrapolation, cross-corpus
robustness, or modality generality.

\section{Conclusion}

Repeatedly exposing a fixed sinusoidal carrier only to query and key
projections is a small, kernel-compatible complement to RoPE. Across three
training seeds it improves language-model NLL throughout the observed context,
and a calibrated low-rank readout adds a further gain beyond a
parameter-matched non-positional control. Component and location ablations
show that RoPE and repeated attention-local access both matter, while exact
logit measurements reveal large content--position interactions. The local
design search is now closed. The next tests should ask whether the frozen
methods survive a new corpus, a modern decoder backbone, and eventually
two-dimensional spatial attention.

\bibliography{references}
\bibliographystyle{iclr2025_conference}

\end{document}
